\documentclass[11pt,a4paper]{article}
\usepackage[margin=23mm]{geometry}
\usepackage{amsmath,amssymb,amsthm,booktabs,microtype,needspace}
\usepackage[T1]{fontenc}
\usepackage{lmodern}
\usepackage[colorlinks=true,linkcolor=blue,urlcolor=blue]{hyperref}
\hypersetup{pdftitle={Subset sums: transported hole bounds and a proof of N at most 96},pdfauthor={Sum-of-Subset research project}}
\newcommand{\FS}{\operatorname{FS}}
\newcommand{\T}{\mathcal T}
\newtheorem{theorem}{Theorem}[section]
\newtheorem{lemma}[theorem]{Lemma}
\newtheorem{proposition}[theorem]{Proposition}
\title{Subset Sums in a Fourfold Sumset\\\large Transported hole bounds and a proof of $N\le96$}
\author{Sum-of-Subset research project}
\date{16 October 2026}
\begin{document}
\maketitle
\begin{abstract}
We exclude cardinality 97 and extend the independently audited chain to
$71\le N\le96$. The proof retains the coarse counting spine of the
preceding simplification, transports positive low holes to 14, and merges
412 retained oriented high-hole states into 17 endpoint groups.
The reflection in the active-center estimate is fixed by an actual
contributor; choosing the numerically stronger orientation is invalid.
Only 12 structural leaf groups need nested envelopes. Seventeen integer
instances require 35 proved binary disjunctions, each with at least one
immediately contradictory child. The resulting 23,303 integer leaf
systems are certified by 156 fresh integer vectors: 2,085 symbolic
interval checks and 52 tree-leaf checks. Every integer model was also
checked separately. The stronger archived bound 84 remains outside the
complete independent audit; the present result advances the audited
chain rather than surpassing that archived bound.
\end{abstract}

\section{Statement, foundation, and what changes}
For a finite real set $A$, let
$\FS(A)=\{\sum_{x\in X}x:X\subseteq A\}$, including zero from the empty
subset, and let $4A=A+A+A+A$ allow repeated summands.
Let $N$ be the largest possible cardinality under $\FS(A)\subseteq4A$.
\begin{theorem}\label{thm:main}
Every such set has at most 96 elements. Hence $71\le N\le96$.
\end{theorem}
The lower bound is the unchanged construction
$[-31,31]\cap\mathbb Z\cup\{\pm93,\pm155,\pm217,\pm713\}$.
The upper-bound foundation is the independently audited chain
$115\to103\to100\to99\to98\to97$, with the latest endpoint and its
simplification in \cite{previous,method}. We reuse its integralization,
covering, peeling, inverse, nonprogression, effective-hole, and span lemmas.
The archived report \cite{archive} provides the strategy of the three
disjunctions below. We reconstruct their validity and generate fresh
certificates; no archived endpoint certificate is read.

A hypothetical size-97 set contains zero: adjoining zero otherwise gives
a forbidden feasible size-98 set. This argument uses the preceding global
upper bound; no hereditary feasibility assertion for arbitrary subsets
is used. Integralize, remove the four outer elements on each side to
obtain $D$, and peel twice to a core
$E=P\cup\{0\}\cup(-Q)$. Write
\begin{gather*}
 d=|D|=89,\quad e=|E|=81,\quad m=a+b=80,\quad
 a=|P|,\ b=|Q|,\quad I=(a-40)^2,\\
 \Lambda=1+T_a+T_b=1641+I,\quad
 \Omega=22d+97-\Lambda=414-I,\quad G=2d-12=166,
\end{gather*}
where $T_r=r(r+1)/2$.
At this endpoint two empty covering units can leave positive low-hole
loss 46, larger than the balanced ordinary-quantile threshold 45.
Thus the previous ordinary-quantile shortcut must be replaced by core
quantiles and, where necessary, proved disjunctions.

\section{Lattice and scalar reduction}
For a positive $r$-element set $X$, retain the definitions
\[
 \mu(X)=|\FS(X)|-1-T_r,\qquad
 \rho(X)=\mu(X)+2-2\delta_0(X),
\]
where $\delta_0(X)$ counts those $x\in X$ for which $2x$ is absent from
the restricted two-sum set of $X\cup\{0\}$.
Set $\rho=\rho(P)+\rho(Q)$, $\mu=\mu(P)+\mu(Q)$, and
$\kappa=|\FS(E)|-\Lambda-\mu$. The inherited nonnegative budget is
\begin{equation}\label{eq:budget}
 \rho+2\theta+c+\T=\Omega,\qquad
 c=2|2D|+18d+85-|\FS(D)|,\qquad
 \T=\kappa+\ell_1+\ell_2+g.
\end{equation}
The one-sided bootstrap is impossible since $1+T_{81}>24\cdot97+77$.
The paired bootstrap gives $q^2\le2288$, with even $q$, hence $q\le46$
and $a,b\ge17$. Nonnegative $\Omega$ leaves 41 oriented pairs,
$a=20,\ldots,60$. For all of them,
$T_{\min(a,b)}-2-\Omega>0$, with minimum 131. The simultaneous inverse
lemma puts $P,Q$ on a common integer lattice, with side support holes
$h_P,h_Q$ satisfying
\begin{equation}\label{eq:holes}
 h_P,h_Q\le\kappa,\quad z=h_P+h_Q-\kappa,
 \quad0\le z\le\min(h_P,h_Q)\le\kappa,
 \quad h_X<|\FS(X)|-1.
\end{equation}

The universal span lemma from \cite{previous} gives
$\max X\le |X|+\rho(X)+h_X+2$ for $|X|\ge15$.
Thus, with $R_p=\max P$ and $R_q=\max Q$,
$R_p+R_q\le m+\rho+\kappa+z+4$.
Before asserting integrality of the peeled layers, the inner comparison is
\[
 |\FS(D_1)|-2|\FS(E)|
 \le C+2\rho-\mu+\kappa+2z+\ell_2,
 \qquad C=2m+4e+14-\Lambda=-1143-I.
\]
For $z=0$, its right side is at most
$C+2\Omega=-315-3I<0$.
For $z>0$, both side supports are nonprogressions by \eqref{eq:holes},
so $\rho\ge52$, $\mu\ge74$. The right side is at most
$C+3\Omega-\rho-\mu\le-27-4I<0$.
An off-lattice point would give two disjoint core-support translates;
this contradiction integralizes the inner peeled layer. The outer-layer
margin is
\[
 \Lambda-(4d-10+\Omega)=881+2I>0.
\]
The inherited overlap and defect lemmas now apply to the integral set $D$:
\begin{equation}\label{eq:propagate}
 w\le4z,\quad \T\ge\kappa+166-3z,\quad
 \T\ge\max\{56,z,166-2z\}.
\end{equation}
Here $w$ counts the support holes of $D$.
If $j$ covering units are empty, put $\delta=c-89j$.
Combining the span lemma with \eqref{eq:budget}--\eqref{eq:propagate} gives
\begin{equation}\label{eq:span}
 R_p+R_q+\delta\le332-I-89j+4z.
\end{equation}

For low holes the scalar budgets are
\begin{equation}\label{eq:low}
 c\le\begin{cases}248-I,&z=0,\\196-I+2z,&1\le z\le14.\end{cases}
\end{equation}
Their maxima are 248 and 224, so $j\le2$.
With one empty unit, $\delta\le159$ or $135$, both below $2\cdot85$.
At most one retained center can therefore have loss at least 85.
With two empty units, $\delta\le70$ or $46$, both below 85, so every
retained center has more than four hull points. The scalar checker
covers all 437 surviving oriented low-hole states.

\section{Transport and high-hole grouping}
\subsection{One monotone change of variables}
Let $M_p=\max D$, $M_q=-\min D$, and let $S_p,S_q$ be its positive
and negative-magnitude masses. For a retained positive center $v$,
$l_v,u_v$ count points of $v+D$ below 1 and above $S_p$; reflection
defines the negative-side quantities. Always
\[
 \sum_v(l_v+u_v)\le\delta,\qquad 0\le l_v+u_v\le88.
\]
Let $K=|2D|$. The cardinality identity gives
\begin{equation}\label{eq:card}
 S_p+S_q+\delta-2K\le1686-89j+4z.
\end{equation}
The mass bonuses are
$\varepsilon_p=\sum P-T_a=\mu(P)+h_P$, and similarly for $q$.

\begin{lemma}[Endpoint transport]\label{lem:transport}
Suppose a group of states has common lower bounds
$\mu(P)\ge u_p$, $\mu(Q)\ge u_q$ and $z\le Z$.
Replace
\[
 S_p\longmapsto\widehat S_p=S_p+Z-z,\qquad
 S_q\longmapsto\widehat S_q=S_q+Z-z.
\]
Its necessary linear rows remain valid at the endpoint $Z$, with mass
bonuses $u_p+Z,u_q+Z$, a common upper bound on the group's loss budget,
and the span and cardinality right sides evaluated at $Z$.
\end{lemma}
\begin{proof}
Put $t=Z-z\ge0$. In a doubled mass row
$2M_p+6R_p-2S_p+2z$, or an exposed mass row
$\sum_{i=1}^4x_i-S_p+z$, the change in the left side is zero.
Other mass and upper quantile left sides decrease by $t$.
The left side of \eqref{eq:card}, after moving $4z$ to the left,
decreases by $2t$; the span expression decreases by $4t$.
For positive low holes the loss-cap expression $\delta-2z$ decreases
by $2t$. For a high-hole group use the maximum of its scalar caps.
All orders, envelope rows, lower quantiles, and losses are unchanged.
Nonnegativity is preserved. These auxiliary masses need not arise from
another actual set; they provide a necessary relaxation.
\end{proof}
For $1\le z\le14$, use $u_p=a-3$, $u_q=b-3$ and $Z=14$.
Then $\varepsilon_p\ge a+11$, $\varepsilon_q\ge b+11$.
The delivered low-hole models have no hole coordinate. With
$B_0=248-89j-I$, their caps and spans become
\[
 \begin{array}{c|cc}
 &\delta\text{ upper bound}&R_p+R_q+\delta\text{ upper bound}\\\hline
 Z&B_0&84+B_0\\
 P&B_0-24&140+B_0
 \end{array}
\]
The $j=2$ core condition is selected from the original budget before
transport; its validity does not depend on interpreting the transported
mass as an actual sum.

\subsection{Effective-hole bounds and the fixed orientation}
For high holes, use the effective-hole witnesses already proved in
\cite{previous}. For clarity, the exact witness test is as follows.
Given $(r,C_*,\gamma)$, set
\begin{gather*}
 K_* =\max\{k\ge3:k(k-3)/2\le C_*+1\},\quad
 D_*=C_*+K_*,\quad E_*=D_*+1,\\
 M_* =\max\left\{2\gamma,1+\left\lfloor
 \frac{1+\sqrt{9+8(E_*+6\gamma)}}2\right\rfloor\right\},
 \quad J_* =\max\{K_*+1,2\gamma+5,M_*\}.
\end{gather*}
The three strict tests are
\[
 T_{M_*}-6\gamma>M_*+1+E_*,\quad
 \gamma(r-J_*+1)>D_*,\quad T_{J_*-1}-1>D_*.
\]
A valid witness gives $\mu(X)\ge r+C_*+1$ when
$2\gamma<h_X<|\FS(X)|-1$.
Search $20\le r\le64$, $7\le\gamma\le20$, $0\le C_*\le600$,
then for each $(r,z)$ choose the largest valid $C_*$ with $2\gamma<z$,
breaking ties by smaller $\gamma$. Each used witness is checked again.
This finite search makes no optimality claim about the witness parameters.

Write $u_r=r+C_*+1$. The inherited flexible inverse, for $2\le\beta\le10$,
is
\[
 v=\max(0,u_r-4\beta+2)=(\beta+1)q+s,
 \quad L_\beta(r,u_r)=\min\{(\beta-1)(r-14),(\beta-1)q+\min(s,\beta-1)\}.
\]
Thus set
$\rho_r^{\rm low}=\max\{r-14,u_r+2-2r,L_2(r,u_r),\ldots,L_{10}(r,u_r)\}$.
Every high state satisfies
\begin{equation}\label{eq:C}
 c\le C_a(z):=414-I-\rho_a^{\rm low}-\rho_b^{\rm low}
                         -\max\{56,z,166-2z\}.
\end{equation}
The loop $15\le z\le414-I$ is exhaustive because $z\le\T\le\Omega$.
It leaves 2,200 nonnegative-budget oriented states, with $C_a(z)\le118<178$.
At most one covering unit is empty.

For an all-active configuration, reflect so that the first paired unit's
actual contributor is positive, and \emph{then} define $a,b$.
The inherited proof gives one of the two lower bounds
\begin{align*}
 F&\ge F_A(a,b)+\tfrac12(u_a+z)+u_b+z,\\
 F&\ge F_B(a,b)+\tfrac23(u_a+z)+u_b+z,
\end{align*}
where
\begin{gather*}
 F=S_p+S_q+c-2(M_p+M_q+R_p+R_q)\le1728+4z,\\
 F_A=\tfrac12T_a+T_b+6a+8b+85,\qquad
 F_B=\tfrac23T_a+T_b+6a+(22b+236)/3.
\end{gather*}
Use the minimum of these two differences from $1728+4z$ in that fixed
orientation. The proof of the active lemma uses nonemptiness to obtain
the contributing slots and outer integrality; the hypothesis $c<d$ in
its earlier application only served to force this nonemptiness.
There is no further need for $c<d$ once all eighteen units are assumed active.

The reflection is not free after $a$ has been assigned. In particular,
at $(a,b,z)=(47,33,112)$ the valid minimum margin is $-1$, while the
reflected minimum is 377. Taking the larger value would incorrectly
remove a necessary model. The checker includes this exact counterexample
as a negative control.
In the fixed contributor orientation, 117 H0 states remain, with
$a=40,\ldots,47$. There are 295 H1 states with $C_a(z)\ge89$;
for them $\delta\le29<\min(a,b)+5$, so ordinary quantiles are valid.

\subsection{Seventeen endpoint groups replace 412 states}
Group states with the same $(a,j,u_a,u_b)$, let $Z$ be the largest
hole count in a group, and let $C^{\max}$ be the largest $C_a(z)$ there.
Lemma~\ref{lem:transport} gives one model with
\[
 \delta\le C^{\max}-89j,\quad
 (\varepsilon_p,\varepsilon_q)\ge(u_a+Z,u_b+Z),\quad
 R_p+R_q+\delta\le332-I-89j+4Z.
\]
Keep \eqref{eq:card} at $Z$. The complete endpoint table is:
\begin{center}\small
\begin{tabular}{rrrrrrr}
\toprule
$a$&$u_a$&$u_b$&H0: $Z$&H0: $C^{\max}$&H1: $Z$&H1: $C^{\max}$\\
\midrule
36&148&224&--&--&64&97\\
37&156&213&--&--&74&107\\
38&165&203&--&--&80&113\\
39&174&193&--&--&84&117\\
40&184&184&174&4&85&118\\
41&193&174&173&14&84&117\\
42&203&165&169&20&80&113\\
43&213&156&163&23&74&107\\
44&224&148&153&21&64&97\\
45&234&139&142&17&--&--\\
46&245&131&128&10&--&--\\
47&256&123&112&0&--&--\\
\bottomrule
\end{tabular}
\end{center}
The scalar enumeration independently reconstructs membership of every
group. A further exact row check verifies transport from every one of
the 412 retained states to its endpoint. The row types are independent
of the outer-center coefficients, so representative modes include both
signs and every relevant row type.
With four zero-empty modes and seventy one-empty modes, the high model
count falls from $117\cdot4+295\cdot70=21118$ to
$8\cdot4+9\cdot70=662$. All 662 endpoint models are excluded by the
coarse model below; none needs a nested envelope or binary tree.

\section{Models and exhaustive selective refinement}
\subsection{Covering modes}
For the outer magnitudes $p_1>\cdots>p_4>M_p$ and their negative
reflection $q_1>\cdots>q_4>M_q$, define
\[
 \lambda_p=p_1-p_2-p_3-p_4,\quad
 \alpha_p=p_2-p_3-p_4,\quad\beta_p=2p_2-p_1-p_3-p_4.
\]
The positive center list is
\[
 v^p=(\lambda_p+p_1,\ldots,\lambda_p+p_4,
       \lambda_p-q_1,\ldots,\lambda_p-q_4,\alpha_p,\beta_p),
\]
with a reflected negative list. String labels \texttt{p0},\ldots,
\texttt{q9} denote centers, not magnitude variables.
Pair 1 is \texttt{p4,q4}. In branch A, $\beta_p+\beta_q\le0$ and
pair 2 is \texttt{p9,q9}; in branch B, $\beta_p+\beta_q>0$ and pair 2
is \texttt{p5,q5}. The other sixteen covering units are singles.

Enumerate every $j$-element set of missing units. If pair 1 is present,
reflect to choose its positive contributor. If pair 2 is present and
nonsplitting, enumerate both contributor choices. In branch A, a missing
single \texttt{p5} or \texttt{q5} can permit pair 2 to split. Drop its
individual losses, enumerate an \emph{actually contributing} member
$s=p,q$, and retain the necessary row $-\beta_s-M_s\le-1$.
If pair 2 itself is missing there is no such member choice.
The unchanged pairing argument in \cite{previous} proves completeness;
there are 4, 70, 578 modes for $j=0,1,2$.
The necessary row for a chosen member is never used as a sufficient
condition for contribution.

The low root families are:
\begin{center}\small
\begin{tabular}{lrrrr}
\toprule
Family&$j$&Integer $a$ range&Groups&Integer models\\
\midrule
Z0&0&25--55&4&124\\
P0&0&26--54&4&116\\
Z1&1&28--52&70&1750\\
P1&1&29--51&70&1610\\
Z2&2&32--48&578&9826\\
P2&2&34--46&578&7514\\
\bottomrule
\end{tabular}
\end{center}
Together with H0 and H1 this gives 1,966 root groups and 21,602 integer
root models. The checker constructs this universe independently of the
certificate routing data.

\subsection{Coarse and nested rows}
All variables are nonnegative. Outer gaps are at least one for $j\le1$,
by the inherited contributing-center integrality argument. For $j=2$
use only weak outer order. Always $M_s\ge R_s+4$, $R_p\ge a$, $R_q\ge b$.
The branch-B row is $-\beta_p-\beta_q\le-1$ for $j\le1$, and its weak
relaxation $\le0$ for $j=2$.
The mass rows, on each side, are
\begin{align}
 S_p&\ge T_a+\varepsilon_p+3R_p+6+M_p,\notag\\
 S_p&\ge T_{a-1}+4R_p+6+M_p.\label{eq:mass}
\end{align}
The first row is multiplied by two in the implementation.
Bonuses are zero for Z, $(a+11,b+11)$ for P, and $(u_a+Z,u_b+Z)$ for H.

Every pair in $2D$, except pairs of two positive or two negative
outliers, lies in $[-M_q-R_q,M_p+R_p]$. There are at most ten unordered
repeated outlier pairs of each sign. Thus
\begin{equation}\label{eq:coarse}
 K\le M_p+M_q+R_p+R_q+21.
\end{equation}
Twice this row added to \eqref{eq:card} eliminates $K$:
\begin{equation}\label{eq:spine}
 S_p+S_q+\delta-2(M_p+M_q+R_p+R_q)
 \le1728-89j+4Z,
\end{equation}
where $Z=0$ in Z, $Z=14$ in P, and $Z$ is the high endpoint in H.
The loss sum, caps, joint span, orders, mass rows and quantiles complete
the coarse model. Unlike the previous simplification, the loss cap and
span rows are retained: some certificates at this endpoint use them.

For a positive center the necessary endpoint rows are:
\begin{center}\small
\begin{tabular}{lll}
\toprule
Type&Lower&Upper\\
\midrule
Global&$-v-l_v-M_p\le-89$&$v-u_v-S_p-M_q\le-88$\\
Core&$-v-l_v-R_p\le-85$&$v-u_v-S_p-R_q\le-84$\\
Ordinary&$-v-l_v\le-b-5$&$v-u_v-S_p\le-a-4$\\
\bottomrule
\end{tabular}
\end{center}
Global rows hold for every nonempty center by integer interval counting.
Core rows hold if its hull count exceeds four, since only four elements
on the relevant side are outside the core; use the core-gap count from
\cite{previous}. Ordinary rows apply at the relevant in-hull endpoint.
Use the reflected rows for negative centers. Individual slot caps are
$\lambda_v\le84$ in a core case, and $\lambda_v\le88$ otherwise.

Initially use global quantiles, except for H1 where ordinary rows have
already been justified. If a two-empty root needs refinement, use core
quantiles throughout, including P2. A one-empty low root can be refined
exhaustively into: no loss $\ge85$, or each retained center in turn is
the unique large-loss center. In the former case use core rows everywhere;
in the latter impose $\lambda_v\ge85$, use global rows there, and core
rows elsewhere. The large-center ranges are ZL: $32\le a\le48$ and
PL: $33\le a\le47$, from \eqref{eq:low} and $\delta\ge85$.

Only six roots need that split. They are Z1 branch A with pair 1 missing
and either pair-2 contributor chosen, Z1/P1 branch A with pair 2 missing,
and Z1/P1 branch A with single \texttt{p0} missing and the negative
pair-2 contributor chosen. Each has 17 retained centers and 18 outcomes.
All other one-empty roots are excluded without naming a large-loss center.

Only twelve resulting structural groups need the nested model.
Expose $x_0=R_p<x_1<x_2<x_3<x_4=M_p$ and its negative analogue $y_i$,
with integral gaps at least one, and add the mass rows
\[
 S_p\ge T_a+\varepsilon_p+\sum_{i=1}^4x_i,\qquad
 S_p\ge T_{a-1}+\sum_{i=0}^4x_i,
\]
with their reflections. For $t=1,2,3$, set $W_t=x_t+y_t$ and $s_t=2(4-t)$.
The inherited interval-envelope proof gives
\[
 K\le2W_t+1+\binom{s_t+1}{2}
 +\sum_{i=t+1}^4\left[
 \min(W_t+1,x_i-x_{i-1})+\min(W_t+1,y_i-y_{i-1})\right].
\]
Indeed the inner hull's double sum has $2W_t+1$ positions, successive
external translates add at most their gap or the hull length, and pairs
among the remaining external points account for the binomial term.
Linearize each minimum by a nonnegative variable bounded above by both
arguments; an actual set chooses the minimum itself. Keep $K$,
\eqref{eq:coarse}, and \eqref{eq:card} in this model. Levels 0 and 4 are
unnecessary. Two extra nonnegative loss coordinates are reserved for
the sole-member disjunction below.

\section{Three disjunctions, including transported upper endpoints}
The disjunctions are necessary for actual configurations. They need not
partition every point of a relaxed real polyhedron. Each child is checked;
no child is discarded on the basis of a numerical solver status.
Every retained contributing center is integral: a contributing point
belongs both to a translate of integral $D$ and to the integral target
support. The original masses, $z$, and the chosen endpoints $Z$ are
integers, so the transported masses are integers as well. These facts
justify the unit gaps below for actual configurations, even though the
linear relaxations allow real coordinates.

\begin{lemma}[Ordinary endpoint or outside endpoint]\label{lem:ordinary}
For a retained positive center $v$, with $S_p$ either the actual or a
transported larger mass, the following disjunctions hold:
\begin{align*}
 -v-l_v\le-b-5\quad&\text{or}\quad(l_v\ge b+5\ \text{and}\ v\le-1),\\
 v-u_v-S_p\le-a-4\quad&\text{or}\quad(u_v\ge a+5\ \text{and}\ v\ge S_p+1).
\end{align*}
Their reflections hold on the negative side.
\end{lemma}
\begin{proof}
If $v\ge0$, at most $v$ of the $b+5$ nonpositive elements of $D$ can
remain above the lower endpoint, giving the first ordinary row.
Otherwise the integral center satisfies $v\le-1$ and all these elements
are lost. At the upper end, when $v\le S_p$, at most $S_p-v$ positive
integers can fit below the endpoint, so the upper ordinary row holds.
When $v>S_p$, integrality gives $v\ge S_p+1$; zero and all $a+4$
positive points are lost.
For transported masses, interpret the count in the enlarged hull.
The recorded upper loss from the original hull is at least the new
upper loss, so the ordinary bound remains valid. If $v>S_p$ for the
enlarged hull, it is also beyond the original endpoint, so at least
$a+5$ original points are lost. One must select the alternative using
the enlarged endpoint, not carry an old outside assertion through the
transport without checking it.
\end{proof}

\begin{lemma}[Inner extreme or endpoint loss]\label{lem:inner}
For $i=1,2,3,4$, the same center satisfies
\[
 l_v\ge i\quad\text{or}\quad v\ge y_{5-i}+1,
 \qquad
 u_v\ge i\quad\text{or}\quad v\le S_p-x_{5-i}.
\]
\end{lemma}
\begin{proof}
If fewer than $i$ points are lost at the lower end, at least one of the
$i$ largest negative magnitudes lies above it. The smallest of those
$i$ is $y_{5-i}$, yielding $v-y_{5-i}\ge1$.
If fewer than $i$ are lost above the original upper endpoint, at least
one of the $i$ largest positive points lies below it, giving
$v+x_{5-i}\le S_p^{\rm original}\le S_p$.
The loss counts are integers, so these two alternatives are exhaustive.
\end{proof}

\begin{lemma}[Both members or the selected sole member]\label{lem:pair}
In a two-empty mode with a potentially split pair 2, let $s$ be the
member chosen from the actual contributors, and let $t$ be the other
side. Either $\beta_t+M_t\ge1$, or there exist hull losses $l_*,u_*\ge0$
for the selected member satisfying
\begin{gather*}
 \sum_{v\text{ retained}}(l_v+u_v)+l_*+u_*\le\delta,\\
 -\beta_s-l_*-R_s\le-85,\qquad
 \beta_s-u_*-S_s-R_t\le-84.
\end{gather*}
\end{lemma}
\begin{proof}
If the other member contributes, the first necessary inequality holds.
Otherwise the selected member is the sole contributor. Its unit's
capacity loss is at least its individual hull loss, so the strengthened
loss sum holds. That member has at least $89-\delta>4$ hull points,
since the two-empty budgets are at most 70 and 46. Its core quantiles
follow. Enlarging the upper mass only weakens its upper row. The
selection of $s$ from actual contributors, rather than from a merely
necessary inequality, is essential to the sole-member conclusion.
\end{proof}

All seventeen integer instances that use a tree appear below.
A range in a cell means each integer separately; the displayed node
counts correspond to those integers in order. All are in branch A.
A large center is present only in ZL and PL and is always \texttt{p8}.
\begin{center}\small
\begin{tabular}{lllll}
\toprule
Family&Missing units&Choice or actual member&$a$&Internal nodes\\
\midrule
ZL&\texttt{p0}&negative contributor&38,39&1,9\\
PL&\texttt{p0}&negative contributor&37--40&1,1,1,1\\
Z2&\texttt{p0,p8}&negative contributor&38,39,40&1,10,2\\
Z2&\texttt{p0,q5}&negative actual member&40--42&1,1,1\\
P2&\texttt{p0,p8}&negative contributor&37--40&1,1,1,1\\
P2&\texttt{p0,q5}&negative actual member&40&1\\
\bottomrule
\end{tabular}
\end{center}
There are 35 internal nodes: 22 inner-extreme, nine ordinary/outside,
and four both/sole-member splits. Every internal node has at least one
terminal child; the longest remaining chain has ten splits. The 17
trees have $17+35=52$ leaves. This structure is discovered numerically,
but its exhaustiveness and every leaf contradiction are checked exactly.

\section{Exact certificates, validation, and conclusion}
For each structural group, all model coefficients are integers and the
right side is an exact rational quadratic polynomial in $a$.
Every stored row-weight vector $u$ consists of positive integers.
For a model $Ax\le b(a)$, $x\ge0$, the checker requires
\[
 uA\ge0\quad\text{coordinatewise},\qquad
 \max_{a\in[L,U]\cap\mathbb Z}ub(a)<0.
\]
The integer maximum uses endpoints and, for a concave quadratic, the
floor and ceiling of its vertex when in range. Polynomials are built
symbolically rather than inferred from samples. Trees occur only at
explicit integer parameters and use exact integer row combinations.

\begin{center}\small
\begin{tabular}{lr}
\toprule
Proof component&Count\\
\midrule
Root groups / integer root models&1966 / 21602\\
Direct coarse root groups&1900\\
Two-empty coarse / nested root refinements&51 / 9\\
One-empty roots requiring a loss split&6\\
Coarse / nested outcomes of those loss splits&105 / 3\\
Final structural groups / nested groups&2068 / 12\\
Integer parameter models before binary trees&23268\\
Tree roots / internal nodes / tree leaves&17 / 35 / 52\\
Integer leaf systems after all refinements&23303\\
Symbolic interval checks / tree-leaf checks&2085 / 52\\
Fresh integer weight vectors&156\\
Least normalized strict gap&$2/299$\\
\bottomrule
\end{tabular}
\end{center}
The last integer-system count replaces each branched parameter by all
its leaves, hence $23268+35=23303$. It should not be confused with the
number of distinct matrix combinations: interval checking avoids their
pointwise repetition. The separate high-state grouping replaces 21,118
potential high models by 662 without omitting an actual state.

The verifier independently constructs all roots, loss outcomes, and
allowed disjunctions. It rejects a missing root or member choice, omitted
large-loss positions, gaps and overlaps between intervals and tree
parameters, missing children, repeated path splits, unknown disjunctions,
invalid row weights, and unjustified model types. It checks high endpoint
transport on all retained states. Verification uses only the Python
standard library and explicit checks, which remain active under
\texttt{-O}. The optional SciPy generator merely discovers candidates;
only exact integer checks permit a certificate to be saved.

The delivered \texttt{--exhaustive} run rebuilds all 23,268 integer
parameter models, compares their rows with the symbolic formulas,
verifies the assigned certificates and all tree leaves, and checks that
each interval's worst gap matches the pointwise minimum. Its output
agrees with the interval run except for the verification-mode label.
Normal and \texttt{-O} outputs are byte-identical; fresh discovery
reproduces the certificate file byte-for-byte; the preceding simplified
upper-97 result is reproduced unchanged. Fifteen negative controls pass,
including the false reflected-margin shortcut, the invalid P2 ordinary
shortcut, and an attempt to certify size 96.
Actual commands, run times, input hashes, regeneration logs, and complete
results are in \texttt{results/2026-10-16/upper96/}.

\Needspace{13\baselineskip}
\begin{proof}[Proof of Theorem~\ref{thm:main}]
The audited foundation excludes all cardinalities above 97. A hypothetical
size-97 set contains zero and satisfies the lattice and scalar reduction.
After the required contributor normalization, a high all-active state is
either excluded by its valid oriented margin or assigned to an H0 group.
All high one-empty states belong to H1. Low states belong to the six
listed low families. Endpoint transport gives a necessary root model.

A root is either directly contradictory or covered by an exhaustive
loss refinement. Any remaining integer instance follows a satisfied
child of each proved disjunction; the transported endpoint arguments
ensure that such a child exists. Every resulting leaf has a strict
nonnegative integer-row contradiction. No size-97 feasible set exists,
so $N\le96$. The existing 71-element construction supplies the lower bound.
\end{proof}

\paragraph{The next endpoint is not certified.}
For size 96, the balanced scalar values would be
$d=88$, $\Omega=432$, $G=164$, so the zero-hole cover cap is 268.
This reaches $3d=264$, permitting three empty units. With one empty
unit the residual cap 180 also reaches $2(d-4)=168$, permitting two
large-loss centers. The present root universe and unique-large-loss
rule therefore do not directly continue to $N\le95$.
These are necessary-budget obstacles, not feasible constructions.

\Needspace{15\baselineskip}
\paragraph{Reproduction.}
From the repository root, with Python 3.10 or later:
\begin{verbatim}
python -I -S -B verification/upper96.py
python -I -S -B verification/upper96.py --exhaustive
\end{verbatim}
The optional generator requires SciPy:
\begin{verbatim}
python experiments/generate_upper96_certificates.py
\end{verbatim}
The Windows entry \texttt{verify\_upper96\_windows.bat} runs the complete
six-stage audited continuation from 115, using the simplified upper-97
step before this final exclusion. Reports and earlier certificates are
retained unchanged for comparison.

\begin{thebibliography}{9}
\bibitem{previous} Sum-of-Subset research project,
\emph{A unified loss and hole proof from 98 to 97}, 16 October 2026.
\href{https://github.com/iH-3/Sum-of-Subset/blob/main/reports/upper97_from98_20261016.pdf}{Repository report}.
Includes references to the audited preceding structural lemmas.
\bibitem{method} Sum-of-Subset research project,
\emph{A shorter proof method for subset sums}, 16 October 2026.
\href{https://github.com/iH-3/Sum-of-Subset/blob/main/reports/proof_method_20261016.pdf}{Repository report}.
\bibitem{archive} Sum-of-Subset research project,
\emph{Complete upper-86 report}, 5 October 2026, sections on the size-97
exclusion and its disjunctions.
\href{https://github.com/iH-3/Sum-of-Subset/blob/main/reports/upper86_complete_20261005.pdf}{Archived report}.
The present finite certificate data were generated independently.
\end{thebibliography}
\end{document}
